\section{Comparison: \texorpdfstring{$\varepsilon$}{epsilon} vs.\ UQLM}
\label{sec:comparison}

\subsection{Setup}
\label{sec:comp-setup}

We compare $\eps$ against UQLM~\cite{uqlm2024}, the closest published tool
for LLM uncertainty quantification. We use UQLM's \texttt{BlackBoxUQ} scorer
with the \texttt{semantic\_negentropy} metric, which is the library's primary
scorer for general-purpose hallucination detection.

Both systems are evaluated on the same 30-prompt benchmark: 10 LOW prompts
(simple, unambiguous utility functions where $\eps$ should not fire) and 20
HIGH prompts (version-split API calls where $\eps$ should fire). The LOW and
HIGH sets are identical to those used in the calibration benchmark of
Section~\ref{sec:epsilon}.

UQLM generates 5 independent responses per prompt and scores their semantic
consistency via a natural language inference model (DeBERTa-large-MNLI).
A confidence score in $[0, 1]$ is returned, where higher values indicate
greater cross-response consistency. We apply a fire threshold of 0.5:
UQLM ``fires'' when confidence $< 0.5$. The $\eps$ system fires when
$\eps_\text{file} \geq 0.30$ (\textsc{flagged}, \textsc{paused}, or
\textsc{aborted}).

\subsection{Results}
\label{sec:comp-results}

\begin{table}[h]
\centering
\caption{Detection rate and false positive rate on the 30-prompt benchmark
         (20 HIGH, 10 LOW). UQLM fires when confidence $< 0.5$;
         $\eps$ fires when $\eps_{\text{file}} \geq 0.30$.
         min\_logprob at threshold $-2.0$ is the best-performing
         single-token baseline.}
\label{tab:comparison}
\small
\begin{tabular}{lcccp{2.6cm}}
\toprule
\textbf{System} & \textbf{HIGH det.}
                & \textbf{LOW FP}
                & \textbf{API calls}
                & \textbf{Localization} \\
\midrule
$\eps$ system        & \textbf{85\%} (17/20) & 50\% (5/10) & 1 & Token, line, function \\
UQLM                 & 5\%   (1/20)           & 0\%  (0/10) & 5 & None (file-level) \\
min\_logprob ($-2.0$)& 45\%  (9/20)           & 30\% (3/10) & 1 & Single token (unfiltered) \\
\bottomrule
\end{tabular}
\normalsize
\end{table}

\begin{figure}[H]
\centering
\includegraphics[width=0.82\textwidth]{fig1_scatter.png}
\caption{$\varepsilon$ score vs.\ UQLM confidence score for all 30 benchmark prompts.
         LOW prompts (blue circles) cluster at low $\varepsilon$; HIGH prompts
         (red triangles) span $\varepsilon \in [0.31, 0.91]$. UQLM confidence
         is $\geq 0.62$ on 29 of 30 prompts regardless of category, providing
         near-zero discriminative signal. The single UQLM detection (pandas
         DataFrame concat) is marked. Points are jittered slightly to resolve
         overlap at discrete UQLM levels.}
\label{fig:scatter}
\end{figure}

UQLM returned a confidence score $\geq 0.62$ on 29 of 30 prompts,
including 19 of 20 HIGH prompts. The single exception was pandas DataFrame
\texttt{concat} (confidence 0.3978), which crossed the 0.5 fire threshold
and constitutes UQLM's only detection. UQLM achieved 5\% detection on the
HIGH set (1/20) and 0\% false positive rate on the LOW set. The lone
detection corresponds to a domain where the pandas \texttt{concat} syntax
genuinely changed between major versions and multiple valid approaches
occasionally appear across independent samples. Figure~\ref{fig:scatter}
shows the full distribution: $\varepsilon$ separates the two categories
cleanly while UQLM confidence is near-maximal for 19 of 20 HIGH prompts.

Lowering the UQLM fire threshold does not recover detection. At threshold
0.9, UQLM fires on additional LOW prompts before covering the remaining
HIGH prompts --- the sensitivity/specificity trade-off never produces
useful discrimination. The confidence scores for HIGH and LOW prompts are
statistically indistinguishable on this benchmark. Figure~\ref{fig:benchmark-bars}
shows the per-prompt breakdown, confirming that UQLM uncertainty is near
zero for almost every prompt in both categories.

Figure~\ref{fig:roc-sweep} makes this argument visual. The $\eps$ threshold
sweep (left panel) exhibits a wide operating region where detection rate
substantially exceeds the false positive rate; the operating point at
$\eps = 0.30$ achieves 85\% detection at 50\% FP. The UQLM threshold sweep
(right panel) shows detection and false positive rates tracking together
with no usable gap: there is no confidence threshold at which UQLM achieves
meaningful detection without proportional false positives.

\begin{figure}[H]
\centering
\includegraphics[width=\textwidth]{fig5_roc_sweep.png}
\caption{Threshold sweep for $\varepsilon$ (left) and UQLM (right) on
         GPT-4o across the 30-prompt benchmark.
         Left: sweeping the $\varepsilon$ fire threshold shows a wide green
         region where detection $>$ FP; the operating point at $t = 0.30$
         achieves 85\% detection at 50\% FP.
         Right: sweeping the UQLM confidence threshold shows detection and
         FP converging with no usable discrimination window --- there is no
         threshold at which UQLM achieves meaningful detection without
         proportional false positives.}
\label{fig:roc-sweep}
\end{figure}

\begin{figure}[H]
\centering
\includegraphics[width=\textwidth]{fig2_benchmark_bars.png}
\caption{Per-prompt $\varepsilon$ (top) and UQLM uncertainty $1 - \text{confidence}$
         (bottom) across the 30-prompt benchmark.
         Top panel bars are coloured by outcome status
         (green: \textsc{complete}; orange: \textsc{flagged}; red: \textsc{paused}).
         The dashed lines mark the soft (0.30) and hard (0.65) $\varepsilon$ thresholds.
         Bottom panel bars are near-zero for 29 of 30 prompts; the single
         UQLM detection (pandas concat, HIGH) is visible as the only bar
         crossing the 0.5 fire threshold.}
\label{fig:benchmark-bars}
\end{figure}

UQLM's 0\% false positive rate reflects the absence of any signal, not the
presence of high precision --- a system that never fires achieves perfect
FPR trivially. The $\eps$ system's 50\% FPR on LOW prompts is the honest
cost of measuring something real: residual Type~3 implementation-approach
uncertainty on utility function prompts that is substantially rarer in
production API integration tasks.

The comparison also highlights a practical asymmetry. UQLM requires 5
API calls per prompt (5$\times$ cost and latency) and returns a single
file-level confidence score with no localization. The $\eps$ system
requires 1 API call, adds zero latency (log-probabilities are returned
with the generation), and provides token-, line-, and function-level
flags with specific alternatives and probabilities. For the code review
use case, localization is not a convenience --- it is the feature that
makes the output actionable.

\subsection{Why UQLM Produces No Signal}
\label{sec:comp-why}

UQLM's \texttt{semantic\_negentropy} measures the semantic consistency of
$N$ independently sampled responses to the same prompt; uncertainty is
inferred from disagreement across generations. For code generation API
decisions, this inference fails. Consider the
Stripe scenario (Scenario~A): at
the method-selector position, the model assigned probability 0.52 to
\texttt{.Ch} (\texttt{Charge}) and 0.47 to \texttt{.Payment}
(\texttt{PaymentIntent}) --- a near-50/50 token-level split. Yet across 5
independent generations the model selects \texttt{PaymentIntent}
consistently, because a 52\%-peaked distribution samples that peak nearly
every time. UQLM observes high consistency and reports confidence $= 1.0$;
the genuine signal is present only at generation time and is gone from the
committed output.

This is the intra- vs.\ inter-generation distinction established in
Section~\ref{sec:intro}: uncertainty that exists at generation time is gone
from the committed output and invisible to multi-sample methods.
It is not a failure of UQLM's implementation but a structural consequence
of measuring uncertainty at the output level rather than at the generation
level.

\subsection{Simple Logprob Baseline}
\label{sec:comp-logprob}

A natural reviewer objection is: \emph{``Why not just threshold on the
minimum token log-probability?''} We evaluate this directly. The
min\_logprob baseline identifies the token with the lowest raw log-probability
in the generation and fires when that value falls below a threshold $\theta$.
No AST filtering, no declaration classification --- raw logprob only.

\begin{table}[h]
\centering
\caption{Min\_logprob baseline detection across thresholds (GPT-4o, 20 HIGH
         prompts / 10 LOW prompts). Best-threshold result (${-2.0}$) compared
         to $\eps$ in Table~\ref{tab:comparison}.}
\label{tab:logprob-baseline}
\small
\begin{tabular}{lcccc}
\toprule
\textbf{Threshold} & \textbf{gpt-4o det.} & \textbf{gpt-4o FP}
                   & \textbf{gpt-4o-mini det.} & \textbf{gpt-4-turbo det.} \\
\midrule
$< -2.0$ & 45\% & 30\% & 40\% & 55\% \\
$< -3.0$ & 25\% & 10\% & 25\% & 35\% \\
$< -4.0$ & 15\% & 0\%  & 10\% & 25\% \\
$< -5.0$ & 10\% & 0\%  & 5\%  & 20\% \\
\midrule
$\varepsilon$ (any) & \textbf{85\%} & 50\% & \textbf{70\%} & \textbf{70\%} \\
\bottomrule
\end{tabular}
\normalsize
\end{table}

At the best threshold for min\_logprob that still detects anything
useful ($-2.0$), detection tops out at 45--55\% across models. $\eps$
achieves 70--85\% --- a gap of 30--40 percentage points at every
operating point. Tightening the threshold to reduce false positives
collapses detection to 10--25\%; loosening it drives false positives
above 30\%.

The failure mode of min\_logprob is structural. A low raw log-probability
identifies the \emph{least certain} token in the generation, but it does
not distinguish whether that token is:
\begin{itemize}[noitemsep]
  \item a function name choice (cosmetic, always harmless),
  \item an API library selector (consequential, may cause failures), or
  \item a whitespace token in a peculiar formatting context (artifact).
\end{itemize}

$\eps$'s AST-based filtering is precisely what closes this gap. Stripping
cosmetic and artifact tokens before aggregation removes the noise that
causes min\_logprob either to over-fire on harmless decisions or to miss
consequential ones by washing them out with noisier alternatives. The
min\_logprob comparison empirically validates the design: the signal
exists in the raw logprobs, but is inaccessible without filtering.
